\documentclass[../../main.tex]{subfiles}
\begin{document}


{\bf [解]}
与えられた関数
\begin{align}
 F(a)= \int_0^{\pi/2}|\sin x-a\cos x|dx \label{eq:1}
\end{align}
について考える．$F(a)$は明らかに連続関数である．
被積分関数の絶対値の中身を
\begin{align}
 f(x)=\sin x - a \cos x \label{eq:2}
\end{align}
とする．一階微分は
\begin{align}
 f'(x) = \cos x + a\cos x \label{eq:3}
\end{align}
である．

まず，$a \le 0$の時は\cref{eq:2}より$f(x) \ge 0$だから\cref{eq:1}の絶対値ははずせて
\begin{align}
 F(a)
   &=\int_0^{\pi/2}f(x) dx \\
   &=\int_0^{\pi/2}(\sin x - a \cos x) dx \\
   &=1-a
\end{align}
である．
これは $a$ の単調減少関数であるから，最大値を考えるときは連続性も考慮に入れて
\begin{align}
 0\le a \label{eq:4}
\end{align}
として良い．

この時は\cref{eq:3}より$f'(x) \ge 0$となって$f(x)$は単調増加である．
また端点では$f(0)=-a\le0$，$f(\pi/2)=1>0$となるから，中間値の定理より
$[0,\pi/2]$ に $f(\alpha)=0$ なる $\alpha$ がただ1つあって，
\begin{align}
 F(a)
  &=-\int_0^\alpha f(x)dx+\int_\alpha^{\pi/2}f(x)dx \\
  &= -\int_0^\alpha \sin x\,dx+\int_0^\alpha a\cos x\,dx+\int_\alpha^{\pi/2}\sin x\,dx-\int_\alpha^{\pi/2}a\cos x\,dx \label{eq:5}
\end{align}
とかける．$\alpha$も$a$の関数であることに注意して両辺を$a$で微分すると
\begin{align}
\begin{split}
 F'(a) 
   &= -\sin\alpha\cdot\alpha'+a\cos\alpha\cdot\alpha'+\int_0^\alpha \cos x\,dx-\sin\alpha\cdot\alpha'+a\cos\alpha\cdot\alpha'-\int_\alpha^{\pi/2}\cos x\,dx \\
   &=2(a\cos\alpha-\sin\alpha)\alpha'+2\sin\alpha-1 \\
   &= 2\sin\alpha-1 \quad (\because f(\alpha)=0)
\end{split}
\end{align}
となって$F(a)$の増減表は以下のようになる．

\begin{table}[htb]
\centering
\caption{$F(a)$ の増減表 ($a \ge 0$)}
\begin{tabular}{c|ccccc}
$a$ & $0$ & $\cdots$ & $\dfrac{\sqrt{3}}{3}$ & $\cdots$ & $(+\infty)$ \\
\hline
$\alpha$ & $0$ & $\cdots$ & $\dfrac{\pi}{6}$ & $\cdots$ & $\left(\dfrac{\pi}{2}\right)$ \\
\hline
$F'(a)$ & $-1$ & $-$ & $0$ & $+$ & $(1)$ \\
\hline
$F(a)$ & $1$ & $\searrow$ & $\sqrt{3}-1$ & $\nearrow$ & 
\end{tabular}
\end{table}

したがって，$\alpha=\dfrac{\pi}{6}$ で $F(a)$ は最小値をとる．
この時，$f(\alpha)=0$ および\cref{eq:2}より，
\begin{align}
 \dfrac{1}{2} - a \dfrac{\sqrt{3}}{2} = 0 \quad \therefore a = \dfrac{\sqrt{3}}{3} \label{eq:5}
\end{align}
である．

従って\cref{eq:1}から$F\left(\dfrac{\sqrt{3}}{3}\right)$を計算すると
\begin{align}
\begin{split}
 F\left(\dfrac{\sqrt{3}}{3}\right) 
  &= -\int_0^{\pi/6}\left(\sin x-\dfrac{\sqrt{3}}{3}\cos x\right)dx+\int_{\pi/6}^{\pi/2}\left(\sin x-\dfrac{\sqrt{3}}{3}\cos x\right)dx \\
  &= -\left[-\cos x-\dfrac{\sqrt{3}}{3}\sin x\right]_0^{\pi/6}+\left[-\cos x-\dfrac{\sqrt{3}}{3}\sin x\right]_{\pi/6}^{\pi/2} \\
  &= \left(\dfrac{\sqrt{3}}{2}+\dfrac{\sqrt{3}}{6}-1\right)-\left(\dfrac{\sqrt{3}}{3}-\dfrac{\sqrt{3}}{2}-\dfrac{\sqrt{3}}{6}\right) \\
  &= \sqrt{3}-1
\end{split}
\end{align}
となる．


\end{document}
